\section{Requirements and System Concept}
\label{sec:requirements}

\subsection{Use case}
A rack holds $N$ perfused bioreactor chips in series inside an incubator. One
imaging head, tilted so its optical axis is horizontal, translates along a
linear X-rail, stops at each chip, focuses through the chip's optical window,
acquires an image (or Z-stack), and runs on-device inference before moving to
the next chip. Acquisition repeats on a schedule over hours to days.

\subsection{Functional requirements}
\begin{itemize}
\item FR1: image each of $N$ chips through its
  $18\times\SI{8}{\milli\meter}$ window.
\item FR2: autofocus per chip via the Z-focus axis.
\item FR3: translate between chips and stop repeatably at each station.
\item FR4: acquire and store time-stamped images with metadata.
\item FR5: run cell detection on-device and log per-chip results.
\item FR6: operate unattended inside an incubator on a fixed schedule.
\end{itemize}

\subsection{Non-functional requirements}
\begin{itemize}
\item NFR1: total parts cost $< \num{500}$~EUR (RQ1).
\item NFR2: built from open hardware and 3D-printed parts; reproducible via
  Ansible.
\item NFR3: no external compute; inference on the Pi~5 + Hailo-8.
\item NFR4: tolerant of the incubator environment (\SI{37}{\celsius}, high
  humidity). \todo{define environmental qualification.}
\end{itemize}

\subsection{Architecture overview}
The system comprises (i) the imaging head---HQ camera, C-mount optics module,
ZEISS objective and Z-focus body; (ii) the motion system---horizontal tilt
mount, V-slot rail, carriage and rail stepper; (iii) the compute stack---Pi~5,
AI~HAT+ (Hailo-8), RTC; and (iv) the software---OpenFlexure Server~v3, the
Ansible provisioning, the scan/time-lapse orchestration and the CVAT annotation
server on a separate VPS. The high-level block diagram is given in the project
\texttt{README.md}; Sections~\ref{sec:optical}--\ref{sec:software} detail each
subsystem.

\subsection{Key design decisions}
\begin{itemize}
\item \emph{Pi~5 + Hailo-8} rather than a Pi~3/4 without NPU: enables on-device
  inference at the cost of an $\approx\num{150}$~EUR premium
  (Section~\ref{sec:evaluation}, Table~\ref{tab:cost}).
\item \emph{OpenFlexure v3} (FastAPI, Python~$\geq 3.11$) rather than the legacy
  v2 server.
\item \emph{Horizontal tilt + linear rail} rather than an XY flexure stage:
  trades lateral flexure motion for multi-chip serial scanning.
\item \emph{Z-focus-only body}: the XY flexure of the stock OpenFlexure body is
  removed, simplifying the print and freeing the sample plane for the rail.
\end{itemize}
See \texttt{docs/RAIL\_MOUNT\_CONCEPT.md} for the rail concept rationale.
